\documentclass[10pt]{article}
\usepackage[T1]{fontenc}
\usepackage[utf8]{inputenc}
\usepackage[french]{babel}
\usepackage{amsmath,amssymb,mathrsfs,mathpazo}
\usepackage{geometry}
\geometry{a4paper,margin=2cm}
\usepackage{graphicx,array,tabularx,tabularray,colortbl,multirow,multicol,booktabs,enumitem,pifont,ulem,fancybox,diagbox,cancel}
\usepackage{pstricks,pstricks-add,pst-plot,pst-tree,pst-node,pst-text,pst-func,pst-math,pst-bspline,pst-3dplot,pst-eucl,pst-solides3d}
\let\pscalculate\undefined
\usepackage{pst-bezier}
\usepackage{tikz,pgfplots,tkz-tab}
\usetikzlibrary{arrows,arrows.meta,patterns,calc,positioning,trees,decorations.markings,shapes.geometric,shapes.misc,shapes.arrows,intersections,decorations.pathmorphing,decorations.pathreplacing,backgrounds,shadows,babel}
\pgfplotsset{compat=1.18}
\usepgfplotslibrary{fillbetween}
\usepackage[np]{numprint}
\usepackage{listings,makecell,pas-tableur,twoopt}
\let\barre\undefined
\usepackage[pstricks]{tablvar}
\definecolor{darkgreen}{RGB}{0,100,0}
\providecommand{\red}{\color{red}}
\providecommand{\blue}{\color{blue}}
\providecommand{\green}{\color{green}}
\providecommand{\black}{\color{black}}
\providecommand{\white}{\color{white}}
\providecommand{\gray}{\color{gray}}
\providecommand{\cyan}{\color{cyan}}
\providecommand{\magenta}{\color{magenta}}
\providecommand{\yellow}{\color{yellow}}
\providecommand{\violet}{\color{violet}}
\providecommand{\orange}{\color{orange}}
\providecommand{\euro}{\texteuro}
\providecommand{\EUR}{\texteuro}
\providecommand{\ieme}{\textsuperscript{e}}
\providecommand{\degres}{\ensuremath{^\circ}}
\providecommand{\degree}{\ensuremath{^\circ}}
\providecommand{\vv}[1]{\overrightarrow{#1}}
\providecommand{\Coord}[1]{#1}
\providecommand{\decosix}{}
\providecommand{\starredbullet}{\ensuremath{\bullet}}
\pagestyle{empty}
\setlength{\parindent}{0pt}
\usepackage{adjustbox}\DeclareUnicodeCharacter{2236}{:}\DeclareUnicodeCharacter{2190}{\ensuremath{\leftarrow}}
\providecommand{\R}{\mathbb{R}}
\providecommand{\N}{\mathbb{N}}
\providecommand{\Z}{\mathbb{Z}}
\providecommand{\Q}{\mathbb{Q}}
\providecommand{\D}{\mathbb{D}}
\providecommand{\C}{\mathbb{C}}
\providecommand{\e}{{\mathrm{e}}}
\providecommand{\vect}[1]{\overrightarrow{#1}}
\providecommand{\vectt}[1]{\overrightarrow{\mathrm{#1}}}
\providecommand{\vvt}[1]{\overrightarrow{\mathrm{#1}}}
\providecommand{\np}[2][]{{#2~#1}}
\providecommand{\vv}[1]{{\overrightarrow{#1}}}
\renewcommand{\d}{{\mathrm{d}}}
\providecommand{\degres}{{^\circ}}
\providecommand{\degre}{{^\circ}}
\providecommand{\blue}{\color{blue}}
\providecommand{\red}{\color{red}}
\providecommand{\enskip}{\quad}
\providecommand{\llbracket}{\lbrack\!\lbrack}
\providecommand{\rrbracket}{\rbrack\!\rbrack}
\providecommand{\euro}{{\text{€}}}
\providecommand{\no}{{\text{n°}}}
\providecommand{\interval}[2]{{[#1;#2]}}
\providecommand{\intervFO}[2]{{[#1;#2[}}
\providecommand{\dotfill}{\ldots}

\providecommand{\ds}{\displaystyle}
\providecommand{\textding}[1]{#1}
\providecommand{\Oij}{\(\mathrm{O};\vec{\imath},\vec{\jmath}\)}
\providecommand{\Oijk}{\(\mathrm{O};\vec{\imath},\vec{\jmath},\vec{k}\)}
\providecommand{\Ouv}{\(\mathrm{O};\vec{u},\vec{v}\)}
\providecommand{\intervFF}[2]{\ensuremath{[#1;#2]}}
\renewcommand{\euro}{\ifmmode\text{\texteuro}\else\texteuro\fi}
\DeclareUnicodeCharacter{00B0}{\ifmmode^{\circ}\else\textdegree\fi}
\DeclareUnicodeCharacter{20AC}{\ifmmode\text{\texteuro}\else\texteuro\fi}
% Énoncés : Éducation nationale. Transcription et corrigés : APMEP, Denis Vergès et contributeurs. Découpage et classement : CoopMaths, Sébastien Lozano, Éric Elter et contributeurs. Adaptation de la présentation : Maths Explorer.
% CC BY-SA : https://coopmaths.fr/www/about/
% https://coopmaths.fr/alea/static/bac/2026/tex/bac_2026_06_sujet1_metropole_1.tex
\begin{document}
\noindent\textit{APMEP / CoopMaths. CC BY-SA. Présentation adaptée par Maths Explorer.}\par\medskip
Un navire assure la liaison entre deux ports de la Méditerranée.

	Lors d'une traversée, une famille a la possibilité de réserver une cabine ainsi qu'un emplacement pour un véhicule.

	\subsection*{Partie A}

	Parmi l'ensemble des familles effectuant la traversée, on constate que $30\,\%$ réservent un emplacement pour un véhicule et, parmi ces dernières, $80\,\%$ réservent une cabine.

	On sait par ailleurs que $75\,\%$ des familles effectuant la traversée réservent une cabine.

	\medskip

	On choisit au hasard une famille effectuant la traversée, et on considère les évènements suivants :

	\begin{itemize}[label={\textbullet~}]
		\item $V$ : \og  la famille réserve un emplacement pour un véhicule \fg{};

		\item $C$ : \og  la famille réserve une cabine \fg{} .
	\end{itemize}

	Pour un évènement quelconque $E$, on désigne par $\overline{E}$ son évènement contraire et par $P(E)$ sa probabilité.

	\begin{enumerate}
		\item \begin{minipage}[t]{0.5\linewidth}
			\begin{enumerate}
			\item Donner $P(C)$.

			\item Reproduire l'arbre ci-contre et compléter les quatre pointillés.
		\end{enumerate}
		\end{minipage}
		\hfill
		
		\begin{tikzpicture}[xscale=1,yscale=1,baseline={(Raa.base)}]
			
			\tikzstyle{noeud}=[fill=white,circle,inner sep=2pt]
			\tikzstyle{feuille}=[fill=white,circle,inner sep=2pt]
			\tikzstyle{etiquette}=[pos=0.6,fill=white, inner xsep=3pt, inner ysep=1.5pt]
			
			\def\DistanceInterNiveaux{2.5}
			\def\DistanceInterFeuilles{0.7}
			
			\def\NiveauA{(0)*\DistanceInterNiveaux}
			\def\NiveauB{(1)*\DistanceInterNiveaux}
			\def\NiveauC{(2)*\DistanceInterNiveaux}
			\def\InterFeuilles{(-1)*\DistanceInterFeuilles}
			
			\node[noeud] (R) at ({\NiveauA},{(1.5)*\InterFeuilles}) {};
			\node[noeud] (Ra) at ({\NiveauB},{(0.5)*\InterFeuilles}) {$V$};
			\node[feuille] (Raa) at ({\NiveauC},{(0)*\InterFeuilles}) {$C$};
			\node[feuille] (Rab) at ({\NiveauC},{(1)*\InterFeuilles}) {$\overline{C}$};
			\node[noeud] (Rb) at ({\NiveauB},{(2.5)*\InterFeuilles}) {$\overline{V}$};
			\node[feuille] (Rba) at ({\NiveauC},{(2)*\InterFeuilles}) {$C$};
			\node[feuille] (Rbb) at ({\NiveauC},{(3)*\InterFeuilles}) {$\overline{C}$};
			
			\draw (R.east)--(Ra.west) node[etiquette] {$\dots$};
			\draw (Ra.east)--(Raa.west) node[etiquette] {$\dots$};
			\draw (Ra.east)--(Rab.west) node[etiquette] {$\dots$};
			\draw (R.east)--(Rb.west) node[etiquette] {$\dots$};
			\draw (Rb.east)--(Rba.west);
			\draw (Rb.east)--(Rbb.west);
		\end{tikzpicture}














		\hfill~

		\item Calculer la probabilité qu'une famille réserve un emplacement pour un véhicule et une cabine.

		\item Une famille a réservé une cabine.

		Déterminer la probabilité qu'elle réserve un emplacement pour un véhicule.

		\item Déterminer $P_{\overline{V}}(C)$.

		On arrondira le résultat à $10^{-2}$.

		Interpréter le résultat dans le contexte de l'exercice.
	\end{enumerate}

	\subsection*{Partie B}
	Sur ce trajet, la réservation d'une cabine et d'un emplacement pour un véhicule sont facturés en supplément du coût de la traversée.

	Ces suppléments sont d'un montant de :

	\begin{itemize}[label={\textbullet~}]
		\item 100~\euro{} pour une cabine ;
		\item 70~\euro{} pour l'emplacement d'un véhicule.
	\end{itemize}

	On suppose qu'une famille peut réserver au maximum une cabine et au maximum un emplacement pour un véhicule.

	À ces suppléments peuvent s'ajouter le prix payé pour des extras (repas, boissons, etc.).

	On note $X$ la variable aléatoire qui associe, à chaque famille effectuant la traversée, le prix qu'elle paie, en euro, pour les suppléments.

	On donne ci-dessous la loi de probabilité de $X$.

	\medskip

	\begin{tblr}{colspec={c *{4}{X[c]}},hlines, vlines}
	$x_{i}$ & 0 & 70 & 100 & 170 \\
	$P\left(X=x_{i}\right)$ & 0,19 & 0,06 & 0,51 & 0,24 \\
	\end{tblr}

	\medskip

	On note $Y$ la variable aléatoire qui associe, à chaque famille effectuant la traversée, le prix qu'elle paie, en euro, pour les extras.

	On admet que la variable aléatoire $Y$ a pour espérance $\mathrm{E}(Y)=104$ et pour variance $\mathrm{V}(Y)=\np{1686}$.

	\medskip

	On suppose que les variables $X$ et $Y$ sont indépendantes.

	\begin{enumerate}
		\item Justifier que \quad $\mathrm{E}(X)=96$ \quad et que \quad $\mathrm{V}(X)=\np{3114}$.

		\item À titre exceptionnel, la compagnie propose une remise de $40\,\%$ sur les suppléments et les extras.

		On note $Z$ la variable aléatoire qui, à chaque famille effectuant la traversée, associe le montant total pour les suppléments et les extras, en euro, payé par cette famille après réduction.

		\begin{enumerate}
			\item Justifier que \quad $Z = 0,6(X + Y)$.
			\item En déduire que \quad $\mathrm{E}(Z)=120$ \quad et que \quad $\mathrm{V}(Z)=\np{1728}$,

			où $\mathrm{E}(Z)$ est l'espérance de la variable aléatoire $Z$ et $\mathrm{V}(Z)$ sa variance.
		\end{enumerate}

		\item On note $n$ un entier naturel non nul et on choisit au hasard un échantillon de $n$ familles effectuant cette traversée bénéficiant de la réduction exceptionnelle définie en question \textbf{2.}

		On admet que ce choix peut être assimilé à un tirage avec remise.

		On désigne par $Z_{1}$ la variable aléatoire égale au prix total pour les suppléments et les extras payé par la première famille, $Z_{2}$ la variable aléatoire égale au prix total pour les suppléments et les extras payé par la deuxième famille et ainsi de suite, $Z_{n}$ la variable aléatoire égale au prix total pour les suppléments et les extras payé par la $n$-ième famille.

		On considère la variable aléatoire \quad $M_{n}=\dfrac{Z_{1}+Z_{2}+\ldots+Z_{n}}{n}$ \quad donnant le prix total moyen pour les suppléments et les extras, payé par ces familles.

		On admet que les variables $Z_{1}$\,, $Z_{2}$\,, \ldots, $Z_{n}$ sont indépendantes et suivent la même loi de probabilité que la variable aléatoire $Z$.

		\begin{enumerate}
			\item Montrer que l'espérance $\mathrm{E}\big(M_{n}\big)$ de la variable $M_{n}$ est égale à 120, et que sa variance $\mathrm{V}\big(M_{n}\big)$ est égale à $\dfrac{\np{1728}}{n}$.

			\item Déterminer le plus petit entier $n$ pour lequel l'inégalité de Bienaymé-Tchebychev permet d'affirmer que \[\mathrm{P}\big(114<M_{n}<126\big) \geqslant 0,85.\]

			Interpréter le résultat dans le contexte de l'exercice.
		\end{enumerate}
	\end{enumerate}

\end{document}
